\documentclass[10pt,a4paper]{amsart}
\usepackage{amsmath,amssymb,amsthm,graphicx}
% Native formulas rendered with standard assembly mathematics packages.
\title{The coinductive monoid of permutation trees is F-inverse}
\author{Antonio Bucciarelli, Arturo De Faveri, Giulio Manzonetto, Antonino Salibra}
\date{doi:10.4230/LIPIcs.FSCD.2026.9; CC BY4.0 (official published article)}
\begin{document}
\maketitle
\noindent\textit{Editorial note.} Source https://drops.dagstuhl.de/entities/document/10.4230/LIPIcs.FSCD.2026.9. Copyright Antonio Bucciarelli, Arturo De Faveri, Giulio Manzonetto and Antonino Salibra, CC BY4.0. Complete permutation-tree F-inverse construction and all author supporting proofs faithfully transcribed; standard external inverse-semigroup criteria remain cited.\par\medskip
\noindent\textit{Allowed standard prerequisites (editorial).}\begin{itemize}
\item Inverse semigroups, regular monoids with commuting idempotents, natural order and minimum group congruence. Publisher section2 and references31/42, including the explicitly cited inverse-monoid criterion [31,Theorem5.1.1]. All actual PT-specific constructions and proofs are included.
\end{itemize}\medskip
\section*{Author conventions (publisher section 2)}
We denote by $\mathbb N$ the set of natural numbers and by $\mathbb N^+$ the set of positive natural numbers.
For every $n\in\mathbb N$, let $S_n$ be the symmetric group on $n$ elements. Let $\operatorname{Sym}(\mathbb N^+)$ be the symmetric group on $\mathbb N^+$.
For every $n\ge1$, $S_n$ can be embedded in $\operatorname{Sym}(\mathbb N^+)$ as the subgroup of all the $\rho\in\operatorname{Sym}(\mathbb N^+)$ such that $\rho i=i$ for all $i>n$. By abuse of notation, we shall write $S_n$ for the image of this embedding. We denote by $\iota$ the identity permutation in $\operatorname{Sym}(\mathbb N^+)$.
Given a semigroup $S$, we say that $u\in S$ is regular if there is $v\in S$ such that $uvu=u$. A semigroup is called regular if every element is regular.
Moreover, we say that $v$ is an inverse of $u$ if $uvu=u$ and $vuv=v$. Every element in a regular semigroup admits at least on inverse.
\paragraph{Definition 1.} An inverse semigroup is a semigroup $S$ endowed with a unary operation $(-)^*$ that satisfies, for all $u,v\in S$:
\[
(u^*)^*=u,\quad(uv)^*=v^*u^*,\quad uu^*u=u,\quad uu^*vv^*=vv^*uu^*.
\]
An inverse monoid is an inverse semigroup together with a unit for multiplication.
In an inverse semigroup $u^*$ is the unique inverse of $u$.
An idempotent in a semigroup $S$ is $u\in S$ such that $uu=u$. We denote by $E(S)$ the set of idempotents of $S$.
\paragraph{The natural order.} By definition, $u\le v$ if $uu^*v=u$, or, equivalently, if $vu^*u=u$. Actually, one can prove that the following are equivalent:
\begin{enumerate}
\item $u\le v$;
\item $u=ve$, for some idempotent $e$;
\item $u=dv$, for some idempotent $d$.
\end{enumerate}
Therefore, if $d$ is idempotent and $u\le d$, then $u$ is idempotent. Moreover, $(S,\le)$ is a partial order compatible with the structure of inverse semigroup, i.e., if $s\le t$ and $u\le v$, then $s^*\le t^*$ and $su\le tv$.
\paragraph{The minimum group congruence.} Two important relations in the study of an inverse semigroup $S$ are:
\begin{enumerate}
\item the relation of compatibility $\sim$, defined by $t\sim u$ if $t^*u$ and $tu^*$ are idempotents;
\item the relation $\sigma_S$, defined by $t\mathrel{\sigma_S}u$ if there exists $w$ such that $w\le t$ and $w\le u$.
\end{enumerate}
When $t$ and $u$ are compatible, i.e., $t\sim u$, then $t\land u$ exists and $t\land u=tu^*u$. The relation $\sigma_S$ is the smallest congruence such that $S/\sigma_S$ is a group [42, Thm.\ 2.4.1] and will be called minimum group congruence.
\paragraph{Definition 4.} An inverse semigroup $S$ is said to be
\begin{enumerate}
\item $E$-unitary if, whenever $e\in E(S)$ and $e\le v$, it follows that $v\in E(S)$;
\item $F$-inverse if each equivalence class of $\sigma_S$ has a unique maximal element.
\end{enumerate}
The inverse semigroup $S$ is $E$-unitary iff $\sim\,=\sigma_S$ [42, Thm.\ 2.4.6]; if $S$ is $F$-inverse, then it is $E$-unitary [42, Proposition 7.1.3].
\paragraph{Trees.} An ordered tree is a non-empty subset $A$ of $(\mathbb N^+)^*$ such that:
\begin{enumerate}
\item[(i)] $A$ is closed under initial segments;
\item[(ii)] $A$ is gapless: for every $i,j\in\mathbb N^+$ and every $w\in A$, if $i<j$ and $wj\in A$, then $wi\in A$.
\end{enumerate}
The elements of $A$ are called nodes; a node is a leaf if it has no children (immediate successors).
The empty word $\epsilon$ is the unique root, i.e., node without immediate predecessors.
We will be exclusively interested in finitely branching trees, i.e., trees whose nodes have a finite number of children.
We denote by $\operatorname{rd}(A)$ the degree of the root of $A$, i.e., the number of its children.
For each $1\le i\le\operatorname{rd}(A)$ we denote by $A_i:=\{w\in(\mathbb N^+)^*:iw\in A\}$ the subtree of $A$ of root $i$.
Finally, we say that a labelled tree is a function $l:A\to L$ from a tree to a nonempty set $L$ of labels. The tree $A$ is called the skeleton of $l$, and is denoted by $\operatorname{sk}(l)$. The set of finitely branching ordered trees can be defined coinductively. When we deal with infinite trees we will freely use coinductive arguments. We refer the reader to [38] for a gentle introduction to coinductive techniques.

\section*{Permutation trees (publisher section 3)}
\paragraph{Definition 12.} A permutation tree is a labelled ordered tree $t:A\to\operatorname{Sym}(\mathbb N^+)$ such that:
\begin{itemize}
\item if $a\in A$ has exactly $n\ge1$ children, then $t(a)\in S_n$;
\item if $a\in A$ is a leaf, then $t(a):=\iota$, where $\iota$ is the identity permutation in $\operatorname{Sym}(\mathbb N^+)$.
\end{itemize}
We let $\mathsf{PT}$ be the set of permutation trees.
We denote by $\bar\epsilon$ the unique permutation tree with skeleton $\{\epsilon\}$, i.e.\ the tree with one node labelled $\iota$.
We can conveniently write a permutation tree $t$ as $\langle\pi;t_1,\ldots,t_n\rangle$ where $\pi\in S_n$ is the label of the root, $\operatorname{rd}(t)=n$ and $t_1,\ldots,t_n\in\mathsf{PT}$.
Observe that $\mathsf{PT}$ can be defined coinductively as the greatest set $\mathsf{PT}$ such that, if $t\in\mathsf{PT}$, then there is $n\in\mathbb N$, $\pi\in S_n$ and $t_1,\ldots,t_n\in\mathsf{PT}$ such that $t=\langle\pi;t_1,\ldots,t_n\rangle$.
\paragraph{Definition 13.} Let $t,u\in\mathsf{PT}$. We define the product $tu$ by coinduction as follows. First, $t\bar\epsilon=t=\bar\epsilon t$.
If $t=\langle\pi;t_1,\ldots,t_n\rangle$, $u=\langle\rho;u_1,\ldots,u_m\rangle$ and $k=\max\{n,m\}$, then we define:
\[
tu:=\langle\pi\circ\rho;u_1t_{\rho1},\ldots,u_kt_{\rho k}\rangle
\]
where, by convention, $t_j=\bar\epsilon$ if $j>n$, and $u_i=\bar\epsilon$ if $i>m$.
The product operation makes $(\mathsf{PT},\cdot,\bar\epsilon)$ into a monoid. We now show that this monoid is an inverse monoid.
\paragraph{Definition 14.} Given $t\in\mathsf{PT}$, we define by coinduction $t^*$ as follows. First, $\bar\epsilon^*=\bar\epsilon$.
If $t=\langle\pi;t_1,\ldots,t_n\rangle$ with $n\ge1$, then $t^*=\langle\pi^{-1};(t_{\pi^{-1}1})^*,\ldots,(t_{\pi^{-1}n})^*\rangle$.
We denote by $E$ the set $\{t\in\mathsf{PT}:t(a)=\iota\text{ for every }a\in\operatorname{sk}(t)\}$.
Every $t\in E$ is univocally determined by its skeleton and we identify $\operatorname{sk}(t)$ with the permutation tree $t$.
\paragraph{Proposition 15.} The set $E(\mathsf{PT})$ of idempotent elements of $\mathsf{PT}$ coincides with $E$. Moreover, for every $d,e\in E$, $de=d\cup e$; consequently any two idempotents commute.
\begin{proof}
Let $t=\langle\pi;t_1,\ldots,t_n\rangle\in E(\mathsf{PT})$ be idempotent. Then $\pi\circ\pi=\pi$, and this is possible only if $\pi=\iota$.
The claim $t\in E$ follows from $t_i=(tt)_i=t_it_i$.
Conversely, if $d=\langle\iota;d_1,\ldots,d_n\rangle\in E$, then by definition of the product $dd=\langle\iota;d_1d_1,\ldots,d_nd_n\rangle$.
Since by coinduction $d_id_i=d_i$, we have $dd=d$ and $E=E(\mathsf{PT})$.
Finally, if $d,e\in E$, then, again by the coinductive definition of the product, $de=d\cup e$.
\end{proof}
\paragraph{Theorem 16.} The monoid $(\mathsf{PT},\cdot,(-)^*,\bar\epsilon)$ is an $E$-unitary inverse monoid.
\begin{proof}
We prove that $\mathsf{PT}$ is an inverse monoid by applying [31, Theorem 5.1.1], which states that $S$ is an inverse monoid iff $S$ is a regular monoid whose idempotents commute.
For every $t=\langle\pi;t_1,\ldots,t_n\rangle\in\mathsf{PT}$, we prove that $tt^*t=t$: $\pi\circ\pi^{-1}\circ\pi=\pi$ and, coinductively,
\[
(tt^*t)_i=((tt^*)t)_i=t_i(tt^*)_{\pi i}=t_i(t^*)_{\pi i}t_{(\pi^{-1}\circ\pi)i}=t_i(t_i)^*t_i=t_i.
\]
Since for every $t$ there is $t^*$ such that $tt^*t=t$, $\mathsf{PT}$ is regular. We conclude that $\mathsf{PT}$ is an inverse monoid applying Proposition 15.
We now prove that $\mathsf{PT}$ is $E$-unitary.
Let $t=\langle\pi;t_1,\ldots,t_n\rangle$, and $e=\langle\iota;e_1,\ldots,e_m\rangle$ be an idempotent such that $e\le t$.
If $e=\bar\epsilon$, then necessarily $t=\bar\epsilon$, hence $t\in E$.
If $e\ne\bar\epsilon$, then $e=td$ for some $d\in E$. We prove that $t\in E$.
First, we immediately get that $\pi=\iota$ and $n\le m$.
Moreover, for every $1\le i\le n$, we have $e_i=d_it_i$.
Since $d_i$ is idempotent, we get that $e_i\le t_i$, and coinductively, that $t_i\in E$. We conclude that $t\in E$.
\end{proof}
As a consequence of $E$-unitarity, if $t\le u$ and one between $t$ and $u$ is idempotent, then both are idempotent.
\paragraph{Lemma 18.} Let $t=\langle\pi;t_1,\ldots,t_n\rangle$ and $u=\langle\rho;u_1,\ldots,u_m\rangle$. Then $t\le u$ iff
\begin{enumerate}
\item $\pi=\rho$ and $m\le n$;
\item $t_i\le u_i$ for $1\le i\le m$ and $t_i\in E$ for $m<i\le n$.
\end{enumerate}
\begin{proof}
Assume that $t\le u$, meaning that $t=ue$ for some $e\in E$. Then $\pi=\rho$ and $n=\max\{m,\operatorname{rd}(e)\}\ge m$.
Moreover, since $t_i=e_iu_i$, for $1\le i\le m$, and $t_i=e_i$, for $m<i\le n$, the conclusion easily follows.
Conversely, let $e=\langle\iota;e_1,\ldots,e_m,t_{m+1},\ldots,t_n\rangle$, where $e_i$ is such that $t_i=e_iu_i$ with $e_i\in E$ for $1\le i\le m$. Then we have $t=ue$.
\end{proof}
Remark that the natural order is the greatest relation on $\mathsf{PT}$ that satisfies conditions (1) and (2) of Lemma 18.
\paragraph{Definition 19.} Given $t=\langle\pi;t_1,\ldots,t_n\rangle\in\mathsf{PT}$, we define coinductively $t^m$ as follows:
\[
t^m:=\begin{cases}
(\langle\pi;t_1,\ldots,t_{n-1}\rangle)^m,&\text{if }t_n\in E\text{ and }\pi n=n,\\
\langle\pi;(t_1)^m,\ldots,(t_n)^m\rangle,&\text{otherwise.}
\end{cases}
\]
Observe the similarity with the definition of Nakajima tree of a $\lambda$-term given in [7, Definition 2.60]. See also [6, Proposition 10.2.15].
Remark that if $e\in E$, then $e^m=\bar\epsilon$.
For $t\in\mathsf{PT}$, let $t^{\uparrow}:=\{u:t\le u\}$.
\paragraph{Proposition 20.} For every $t,u\in\mathsf{PT}$, the following hold:
\begin{enumerate}\item $t\le t^m$;\item if $t\le u$, then $u\le t^m$.\end{enumerate}
In particular, $t^m$ is the unique maximal element in $t^{\uparrow}$.
\begin{proof}
If $t$ is idempotent, (1) is trivial and (2) follows from $E$-unitarity.
Otherwise let $t=\langle\pi;t_1,\ldots,t_n\rangle$ and let $k$ be the greatest $i$ such that $t_k\notin E$ or $\pi k\ne k$.
Then $t^m=\langle\pi;(t_1)^m,\ldots,(t_k)^m\rangle$.
By coinductive hypothesis $t_i\le(t_i)^m$ for $1\le i\le k$, and $t_i\in E$ for $k<i\le n$, so that by Lemma 18 $t\le t^m$, showing (1).
Concerning (2), let $u=\langle\rho;u_1,\ldots,u_m\rangle$ such that $t\le u$.
Without loss of generality, we assume that $n=m$, so that $k\le m$, where $k$ is defined above.
Moreover, we have $\pi=\rho$, so that by Lemma 18 and by coinductive hypothesis, $u_i\le(t_i)^m$ for $1\le i\le k$.
As $t_i\le u_i$ and $t_i\in E$ for $k<i\le n$, we have $u_i\in E$ for $k<i\le m$ by $E$-unitarity.
By Lemma 18 we conclude that $u\le t^m$.
\end{proof}
We now prove that $\mathsf{PT}$ is $F$-inverse. The proof relies on the following general lemma that will be used multiple times.
\paragraph{Lemma 21.} Let $S$ be a $E$-unitary inverse semigroup. Assume that for every $t\in S$, the set $t^{\uparrow}$ has a unique maximal element $t^m$.
Then for all $t,u\in S$, $t\mathrel{\sigma_S}u$ iff $t^m=u^m$.
\begin{proof}
Assuming $t\mathrel{\sigma_S}u$, there exists $w$ such that $w\le t,u$. Hence $w^m=t^m=u^m$.
Conversely, let $w=t^m=u^m$. Then $u=we$ for some $e\in E(S)$ and $t=dw$ for some $d\in E(S)$.
Moreover, $t^*u=t^*we\le w^*we\in E(S)$ and $tu^*=dwu^*\le dww^*\in E(S)$, so that $t,u$ are compatible, i.e., $t\sim u$.
As $S$ is $E$-unitary, $\sim\,=\sigma_S$, so that $t\mathrel{\sigma_S}u$.
\end{proof}

\paragraph{Theorem 22.} The monoid $\mathsf{PT}$ is $F$-inverse.

\begin{proof}
By Lemma 21 and Proposition 20.
\end{proof}

\begin{thebibliography}{99}
\bibitem[1]{source1} Samson Abramsky. Computational interpretations of linear logic. Theor. Comput. Sci., 111(1\&2):3-57, 1993. URL: https://doi.org/10.1016/0304-3975(93)90181-R .
\bibitem[2]{source2} Roberto Amadio and Pierre-Louis Curien. Domains and lambda-calculi, volume 46 of Cambridge Tracts in Theoretical Computer Science. CUP, 1998.
\bibitem[3]{source3} Andrea Asperti and Giuseppe Longo. Categories, types and structures: an introduction to category theory for the working computer scientist. MIT Press, Cambridge, MA, 1991.
\bibitem[4]{source4} Karl Auinger, Ganna Kudryavtseva, and Mária B. Szendrei. F-inverse monoids as algebraic structures in enriched signature. Indiana University Mathematics Journal, 70(5):2107-2131, 2021.
\bibitem[5]{source5} Hendrik Pieter Barendregt, Wil Dekkers, and Richard Statman. Lambda Calculus with Types. Perspectives in logic. Cambridge University Press, 2013. URL: http://www.cambridge.org/de/academic/subjects/mathematics/logic-categories-and-sets/lambda-calculus-types .
\bibitem[6]{source6} Henk Barendregt. The Lambda Calculus: Its Syntax and Semantics, volume 103 of Studies in Logic and the Foundations of Mathematics. North-Holland, 1984.
\bibitem[7]{source7} Henk Barendregt and Giulio Manzonetto. A Lambda Calculus Satellite, volume 94 of Studies in logic: Mathematical logic and foundations. College Publications, 2022. URL: https://www.collegepublications.co.uk/logic/mlf/?00035 .
\bibitem[8]{source8} Jan A. Bergstra and Jan Willem Klop. Invertible terms in the lambda calculus. Theor. Comput. Sci., 11:19-37, 1980. URL: https://doi.org/10.1016/0304-3975(80)90032-8 .
\bibitem[9]{source9} Lison Blondeau-Patissier, Pierre Clairambault, and Lionel Vaux Auclair. Strategies as resource terms, and their categorical semantics. Log. Methods Comput. Sci., 21(4), 2025. URL: https://doi.org/10.46298/LMCS-21(4:9)2025 .
\bibitem[10]{source10} Flavien Breuvart, Giulio Manzonetto, Andrew Polonsky, and Domenico Ruoppolo. New results on Morris’s observational theory: The benefits of separating the inseparable. In 1st International Conference on Formal Structures for Computation and Deduction, FSCD 2016, volume 52 of Leibniz International Proceedings in Informatics (LIPIcs), pages 15:1-15:18. Schloss Dagstuhl - Leibniz-Zentrum für Informatik, 2016. URL: https://doi.org/10.4230/LIPIcs.FSCD.2016.15 .
\bibitem[11]{source11} Corrado Böhm and Mariangiola Dezani-Ciancaglini. Combinatory logic as monoids. Fundamenta Informaticae, 12(4):525-540, 1989.
\bibitem[12]{source12} Alonzo Church. The Calculi of Lambda Conversion. Princeton University Press, 1941.
\bibitem[13]{source13} Pierre Clairambault. Causal Investigations in Interactive Semantics. HAL, 2024. URL: https://tel.archives-ouvertes.fr/tel-04523273 .
\bibitem[14]{source14} Mario Coppo, Mariangiola Dezani-Ciancaglini, and Maddalena Zacchi. Type theories, normal forms, and D\_∞-lambda-models. Information and Computation, 72:85-116, 1987.
\bibitem[15]{source15} Roberto Di Cosmo. A short survey of isomorphisms of types. Math. Struct. Comput. Sci., 15(5):825-838, 2005. URL: https://doi.org/10.1017/S0960129505004871 .
\bibitem[16]{source16} Lukasz Czajka. A new coinductive confluence proof for infinitary lambda calculus. Log. Methods Comput. Sci., 16(1), 2020. URL: https://doi.org/10.23638/LMCS-16(1:31)2020 .
\bibitem[17]{source17} Daniele D'Angeli, Emanuele Rodaro, and Jan Philipp Wächter. On the structure theory of partial automaton semigroups. Semigroup Forum, 101:51-76, 2020.
\bibitem[18]{source18} Vincent Danos and Laurent Régnier. Local and asynchronous beta-reduction. In Proceedings of the 8th Annual IEEE Symposium on Logic in Computer Science (LICS 1993), pages 296-306, USA, 1993. IEEE Computer Society.
\bibitem[19]{source19} Vincent Danos and Laurent Régnier. Reversible, irreversible and optimal lambda-machines. Theoretical Computer Science, 227(1-2):77-97, 1999.
\bibitem[20]{source20} Mariangiola Dezani-Ciancaglini. Characterization of normal forms possessing inverse in the \ensuremath{\lambda} \ensuremath{\beta} \ensuremath{\eta}-calculus. Theoretical Computer Science, 2(3):323-337, 1976.
\bibitem[21]{source21} Alejandro Díaz-Caro, Mauricio Guillermo, Alexandre Miquel, and Benoît Valiron. Realizability in the unitary sphere. In 34th Annual ACM/IEEE Symposium on Logic in Computer Science, LICS 2019, Vancouver, BC, Canada, June 24-27, 2019, pages 1-13. IEEE, 2019. URL: https://doi.org/10.1109/LICS.2019.8785834 .
\bibitem[22]{source22} Thomas Ehrhard. On Köthe sequence spaces and linear logic. Math. Struct. Comput. Sci., 12(5):579-623, 2002. URL: https://doi.org/10.1017/S0960129502003729 .
\bibitem[23]{source23} Enno Folkerts. Invertibility in \ensuremath{\lambda}\ensuremath{\eta}. In Thirteenth Annual IEEE Symposium on Logic in Computer Science, Indianapolis, Indiana, USA, June 21-24, 1998, pages 418-429. IEEE Computer Society, 1998.
\bibitem[24]{source24} Jonathon Funk. Semigroups and toposes. Semigroup Forum, 75:480-519, 2007.
\bibitem[25]{source25} Jonathon Funk and Benjamin Steinberg. The universal covering of an inverse semigroup. Applied Categorical Structures, 18:135-163, 2010. URL: https://doi.org/10.1007/s10485-008-9147-2 .
\bibitem[26]{source26} Pietro Di Gianantonio, Gianluca Franco, and Furio Honsell. Game semantics for untyped \ensuremath{\lambda} \ensuremath{\beta} \ensuremath{\eta}-calculus. In Jean-Yves Girard, editor, Typed Lambda Calculi and Applications, 4th International Conference, TLCA'99, L'Aquila, Italy, April 7-9, 1999, Proceedings, volume 1581 of Lecture Notes in Computer Science, pages 114-128. Springer, 1999.
\bibitem[27]{source27} Jean-Yves Girard. Linear logic. Theoretical Computer Science, 50(1):1-101, 1987. URL: https://doi.org/10.1016/0304-3975(87)90045-4 .
\bibitem[28]{source28} Jean Goubault-Larrecq. Musings around the geometry of interaction, and coherence. Theoretical Computer Science, 412(20):1998-2014, 2011. URL: https://doi.org/10.1016/j.tcs.2010.12.023 .
\bibitem[29]{source29} Giulio Guerrieri and Federico Olimpieri. Categorifying non-idempotent intersection types. In Christel Baier and Jean Goubault-Larrecq, editors, 29th EACSL Annual Conference on Computer Science Logic, CSL 2021, Ljubljana, Slovenia (Virtual Conference), January 25-28, 2021, volume 183 of LIPIcs, pages 25:1-25:24. Schloss Dagstuhl - Leibniz-Zentrum für Informatik, 2021. URL: https://doi.org/10.4230/LIPIcs.CSL.2021.25 .
\bibitem[30]{source30} Karl H. Hofmann and Michael W. Mislove. All compact Hausdorff lambda models are degenerate. Fundamenta Informaticae, 22(1,2):23-52, 1995. URL: https://doi.org/10.3233/FI-1995-22122 .
\bibitem[31]{source31} John M. Howie. Fundamentals of Semigroup Theory, volume 12 of London Mathematical Society monographs. Oxford Science Publications, 1995.
\bibitem[32]{source32} John Martin E. Hyland. A syntactic characterization of the equality in some models for the \ensuremath{\lambda}-calculus. Journal London Mathematical Society (2), 12(3):361-370, 1975.
\bibitem[33]{source33} Benedetto Intrigila, Giulio Manzonetto, and Andrew Polonsky. Refutation of Sallé’s longstanding conjecture. In 2nd International Conference on Formal Structures for Computation and Deduction, FSCD 2017, volume 84 of Leibniz International Proceedings in Informatics (LIPIcs), pages 20:1-20:18. Schloss Dagstuhl - Leibniz-Zentrum für Informatik, 2017. URL: https://doi.org/10.4230/LIPIcs.FSCD.2017.20 .
\bibitem[34]{source34} Benedetto Intrigila, Giulio Manzonetto, and Andrew Polonsky. Degrees of extensionality in the theory of Böhm trees and Sallé’s conjecture. Logical Methods in Computer Science, 15(1), 2019. URL: https://doi.org/10.23638/LMCS-15(1:6)2019 .
\bibitem[35]{source35} Benedetto Intrigila and Monica Nesi. On structural properties of eta-expansions of identity. Information Processing Letters, 87(6):327-333, 2003. URL: https://doi.org/10.1016/S0020-0190(03)00344-2 .
\bibitem[36]{source36} Peter R. Jones, Stuart W. Margolis, John Meakin, and Joseph B. Stephen. Free products of inverse semigroups II. Glasgow Mathematical Journal, 33(3):373-387, 1991.
\bibitem[37]{source37} Axel Kerinec, Giulio Manzonetto, and Federico Olimpieri. Why are proofs relevant in proof-relevant models? Proc. ACM Program. Lang., 7(POPL):218-248, 2023. URL: https://doi.org/10.1145/3571201 .
\bibitem[38]{source38} Dexter Kozen and Alexandra Silva. Practical coinduction. Mathematical Structures in Computer Science, 27(7):1132-1152, 2017. URL: https://doi.org/10.1017/S0960129515000493 .
\bibitem[39]{source39} Jean-Louis Krivine. On the structure of classical realizability models of ZF. In Hugo Herbelin, Pierre Letouzey, and Matthieu Sozeau, editors, 20th International Conference on Types for Proofs and Programs, TYPES 2014, Paris, France, May 12-15, 2014, volume 39 of LIPIcs, pages 146-161. Schloss Dagstuhl - Leibniz-Zentrum für Informatik, 2014. URL: https://doi.org/10.4230/LIPIcs.TYPES.2014.146 .
\bibitem[40]{source40} Joachim Lambek and Philip J. Scott. Introduction to higher order categorical logic, volume 007 of Cambridge studies in advanced mathematics. CUP, Cambridge, New York, Melbourne, 1986.
\bibitem[41]{source41} Adrienne Lancelot, Beniamino Accattoli, and Maxime Vemclefs. Barendregt’s theory of the \ensuremath{\lambda}-calculus, refreshed and formalized. In Yannick Forster and Chantal Keller, editors, 16th International Conference on Interactive Theorem Proving, ITP 2025, Reykjavik, Iceland, September 28 - October 1, 2025, volume 352 of LIPIcs, pages 13:1-13:22. Schloss Dagstuhl - Leibniz-Zentrum für Informatik, 2025. URL: https://doi.org/10.4230/LIPIcs.ITP.2025.13 .
\bibitem[42]{source42} Mark V. Lawson. Inverse Semigroups. World Scientific, 1998.
\bibitem[43]{source43} Mark V. Lawson. Non-commutative Stone duality: inverse semigroups, topological groupoids and C*-algebras. International Journal of Algebra and Computation, 22(6):1-47, 2012. URL: https://doi.org/10.1142/S0218196712500580 .
\bibitem[44]{source44} Stefania Lusin and Antonino Salibra. The lattice of lambda theories. J. Log. Comput., 14(3):373-394, 2004. URL: https://doi.org/10.1093/logcom/14.3.373 .
\bibitem[45]{source45} Giulio Manzonetto and Antonino Salibra. Applying universal algebra to lambda calculus. Journal of Logic and Computation, 20(4):877-915, 2010. URL: https://doi.org/10.1093/logcom/exn085 .
\bibitem[46]{source46} James Hiram Morris. Lambda calculus models of programming languages. Phd thesis, Massachusetts Institute of Technology (MIT), 1968.
\bibitem[47]{source47} Simona Ronchi Della Rocca and Luca Paolini. The Parametric Lambda Calculus: A Metamodel for Computation. Texts in Theoretical Computer Science. An EATCS Series. Springer Berlin, Heidelberg, 2004.
\bibitem[48]{source48} Antonino Salibra. On the algebraic models of lambda calculus. Theoretical Computer Science, 249:197-240, 2000. URL: https://doi.org/10.1016/S0304-3975(00)00059-1 .
\bibitem[49]{source49} Rick Statman. Church’s semigroup is Sq-universal. In Proceedings of the 6th International Conference on Formal Structures for Computation and Deduction (FSCD 2021), volume 196 of Leibniz International Proceedings in Informatics (LIPIcs), pages 6:1-6:6. Schloss Dagstuhl - Leibniz-Zentrum für Informatik, 2021. URL: https://doi.org/10.4230/LIPIcs.FSCD.2021.6 .
\bibitem[50]{source50} Steffen van Bakel, Franco Barbanera, Mariangiola Dezani-Ciancaglini, and Fer Jan de Vries. Intersection types for \ensuremath{\lambda}-trees. Theoretical Computer Science, 272:3-40, 2002.
\end{thebibliography}
\end{document}
